\chapter{Introduction}\label{ch:intro}

Many community buildings are no longer just passive electricity users. Modern
homes can have solar panels, batteries, HVAC equipment for heating and cooling,
and other controllable devices that support local energy flexibility.
If homes in a community are grouped and controlled together, their total
electricity use can be moved away from peak periods and made closer to a
desired target load profile. The challenge is that each home only sees its own
situation, but the success of the controller is judged by the behaviour of the
whole building portfolio.

This midterm report studies portfolio load tracking in the CityLearn simulation
environment. CityLearn is an open benchmark for building energy control and for
comparing reinforcement learning algorithms
\cite{citylearn2019,citylearn2020,citylearnv2}. The project uses the IEA EBC
Annex 96 Common Exercise 1 (CE1) as a benchmark setting. In this
task, 25 single-family homes need to follow a daily reference load profile as a
group. The controller can adjust battery operation and HVAC-related flexibility
for both heating and cooling. The evaluation considers not only load tracking
error, but also indoor comfort and secondary metrics such as total energy use,
peak demand, peak-to-valley ratio, load factor, ramping, and fairness. Cost and
carbon emissions are part of the Annex 96 metric set, but the CE1 schemas used
in this project do not provide pricing or carbon-intensity files, so these two
fields are not interpreted in the current experiments.

The core research question is:

\begin{quote}
Can multi-agent reinforcement learning improve the CityLearn CE1 load tracking
task, what type of MARL is most suitable for this problem, and can structured
communication between building agents further improve performance?
\end{quote}

This project is organized around three comparisons. First, it compares
rule-based control and independent building control with MARL-based coordination
to test whether learned coordination is actually useful for CE1. Second, using
the stable on-policy method MAPPO as the current example, it compares different
MARL structures, including grouped and ungrouped policies, communication and
no-communication variants, and intra-group and inter-group communication. Third,
it will compare different types of MARL, such as MAPPO, MASAC, or value-based
methods, to understand which learning style is more suitable for this task.

At the current stage, MAPPO is used as the main experimental backbone. This is
a practical choice rather than a final claim that MAPPO is the best algorithm.
MAPPO is an on-policy centralized-training, decentralized-execution method, and
PPO-style updates are usually stable and easier to debug
\cite{ppo2017,mappo2022}. This is useful for first building a reliable
CityLearn pipeline with grouped actors, a shared critic, and communication
modules. After this pipeline is stable, the project can compare MAPPO with
other directions such as SAC/MASAC or DQN-style value-based methods.

The communication part is therefore an optimization and innovation on top of
the current MAPPO backbone. It studies whether agents should share information,
and if so, whether simple mean-message communication, PowerNet-style structured
information passing, graph-attention communication, attention-based
communication, or weighted same-group/other-group communication is more
effective.

The current stage of the project has produced a working research codebase and a
clear experimental plan. The completed work includes:

\begin{itemize}
  \item an RBC notebook baseline for non-learning comparison in Texas and
  Vermont;
  \item implementation of CityLearn wrappers for independent PPO/SAC and MAPPO
  training;
  \item implementation of standard MAPPO and grouped MAPPO baselines;
  \item implementation of grouped actor policies using K-means clustering over
  building flexibility features;
  \item implementation of several differentiable communication modules,
  including CommNet-style mean pooling, PowerNet-style structured information
  passing, GAT-style graph attention,
  TarMAC-style attention, DIAL-inspired global communication, and weighted
  same-group/other-group communication;
  \item implementation of logging, checkpointing, testing, and Annex 96 metric
  export utilities.
\end{itemize}

The project now has a completed Vermont benchmark comparison for the selected
controllers, together with exported test-month metrics, ranking tables, and load
tracking figures generated from the \texttt{results}, \texttt{wandb}, and
\texttt{target\_folder} outputs. These results move the work beyond pipeline
debugging and allow a more mature comparison of tracking, comfort, peak demand,
fairness, and communication design. The current strongest overall Vermont result
is grouped MAPPO with the \texttt{comm\_v2} communication variant, while TarMAC
and the weighted \mbox{$\alpha=0.90, \beta=0.10$} setting are also competitive.
Texas RL experiments and multi-seed validation remain future extensions, but the
Vermont comparison reported here is now a completed benchmark set rather than a
debugging-stage pilot.

Chapter~\ref{ch:background} reviews the relevant literature and explains where
the project fits. Chapter~\ref{ch:method} describes the implemented method and
current progress. Chapter~\ref{ch:eval} gives the completed Vermont evaluation
and the main result analysis. Chapter~\ref{ch:concl} summarizes the midterm
outcomes and next steps.
